\chapter{Galactosemia Screening}

\section*{What Is This Test?}
Galactosemia screening is a newborn screening test that detects classical galactosemia, an autosomal recessive inborn error of galactose metabolism caused by a severe deficiency of the enzyme galactose-1-phosphate uridylyltransferase (GALT). GALT catalyzes the conversion of galactose-1-phosphate and UDP-glucose to UDP-galactose and glucose-1-phosphate, the second step in the Leloir pathway of galactose metabolism. In GALT deficiency, galactose-1-phosphate accumulates in the liver, brain, kidneys, and the lenses of the eyes, causing: acute neonatal illness (jaundice, vomiting, failure to thrive, hepatomegaly, coagulopathy, Escherichia coli sepsis, and death if untreated within the first 1--2 weeks of life), long-term complications despite early treatment (primary ovarian insufficiency in >80\% of females, developmental delay, speech dyspraxia, and neurologic complications), and bilateral cataracts from the accumulation of galactitol (the reduced product of galactose formed via the aldose reductase pathway) in the lens. Newborn screening is performed on dried blood spots collected by heel stick at 24--48 hours of life. Galactosemia is included in the Recommended Uniform Screening Panel (RUSP) and is screened in all 50 states \cite{kaye2006newborn}. Incidence: approximately 1:30,000 to 1:60,000 in the United States (classical galactosemia); higher in Irish Traveller populations. The treatment is immediate and lifelong elimination of galactose from the diet (no breast milk, no cow's milk formula, no dairy products; soy formula is substituted). Newborn screening allows detection BEFORE the onset of life-threatening neonatal illness.

\section*{Alternative Names}
GALT Deficiency; Classical Galactosemia; Galactose-1-Phosphate Uridylyltransferase Deficiency; Type 1 Galactosemia; Newborn Screening for Galactosemia; Beutler Test (historical); Galactosemia Screen

\section*{Principle}
Newborn screening for galactosemia is performed on dried blood spots. The primary screening method is the measurement of GALT enzyme activity using a fluorometric or colorimetric assay in which GALT activity in the dried blood spot is coupled to an NADH-producing reaction and measured spectrophotometrically or fluorometrically. The Beutler test (historical) used a fluorescent spot test: a dried blood spot is incubated with galactose-1-phosphate, UDP-glucose, NADP, and auxiliary enzymes; GALT activity produces NADPH which fluoresces under UV light. Absence of fluorescence indicates GALT deficiency. Currently, most U.S. newborn screening programs use tandem mass spectrometry (MS/MS) or quantitative fluorometry. A second-tier test measures the total galactose (galactose + galactose-1-phosphate) concentration in the dried blood spot by MS/MS or by enzymatic methods. Confirmatory testing in a positive screen or symptomatic infant is performed by: quantitative GALT enzyme activity in red blood cells (RBCs, the gold standard) --- classical galactosemia shows GALT activity <1\% of normal, Duarte variant shows 15--30\% activity. Molecular genetic testing of the GALT gene (chromosome 9p13, most common pathogenic variant in Caucasians: Q188R, accounting for 60--70\% of classical galactosemia alleles; S135L in African Americans; Duarte variant [N314D] with the Duarte-2 [D2] variant reduces GALT activity but does not cause disease).

\section*{History}
Galactosemia was first described in 1908 by von Reuss as an infant with galactosuria, failure to thrive, and hepatomegaly. In 1917, Goppert described an autosomal recessive inheritance pattern. In 1935, Mason and Turner described the association of galactosemia with cataracts. In 1956, Herman Kalckar and colleagues elucidated the Leloir pathway of galactose metabolism; the enzyme deficiency was identified by Kurt Isselbacher and colleagues in 1956. Development of the galactose-restricted diet for the treatment of galactosemia was described in the 1950s. The Beutler fluorescent spot test for GALT enzyme deficiency was introduced in 1964, enabling newborn screening \cite{beutler1965new}. Universal newborn screening for galactosemia was implemented in the United States in the 1960s--1970s \cite{fridovichkeil2019galactosemia}.

\section*{Why This Test Is Ordered}
Universal newborn screening for galactosemia on all newborns (mandatory). Diagnostic testing for galactosemia in a neonate presenting with jaundice, hepatomegaly, vomiting, coagulopathy, cataracts, E. coli sepsis, or failure to thrive in the first 1--2 weeks of life. Confirmatory testing after a positive newborn screen (quantitative GALT enzyme in RBCs). Monitoring of galactose-1-phosphate (Gal-1-P) levels in RBCs in patients with known galactosemia to assess dietary compliance.

\section*{Primary vs Secondary Test}
This is a primary newborn screening test. Universal screening is mandatory in the United States.

\section*{How This Test Is Performed}
Newborn screening: heel stick blood collected on filter paper (Guthrie card) at 24--48 hours of life, after the infant has received breast milk or formula (galactose-containing feeds). Diagnostic testing: venous blood sample for quantitative GALT enzyme activity in RBCs (sodium heparin, green-top tube, must be sent on ice immediately to the lab). GALT enzyme activity can be falsely low after a packed red blood cell transfusion. For infants who have been transfused, molecular genetic testing of the GALT gene from leukocyte DNA (pre-transfusion) or from cultured fibroblasts is required.

\section*{What Preparation Is Needed}
No special preparation is required for the newborn heel-stick specimen. The infant should have received breast milk or formula (galactose-containing feeds) for at least 24 hours before collection, because screening before adequate feeding can yield falsely normal GALT results. For diagnostic quantitative GALT enzyme testing, no fasting is needed, but the laboratory must be aware of any recent packed red blood cell transfusion (which produces falsely low enzyme activity) and the specimen must be sent on ice immediately.

\section*{Contraindications and Precautions}
\begin{itemize}
  \item There are no contraindications to newborn screening or venipuncture. Important interpretive cautions:
  \item (1) GALT enzyme activity is falsely low after a packed red blood cell transfusion, so molecular genetic testing of the GALT gene is required in transfused infants.
  \item (2) A positive screen is a medical emergency: dietary intervention must not be delayed while awaiting confirmatory results \cite{berry2012galactosemia}.
  \item (3) Not all elevated total galactose reflects GALT deficiency --- galactokinase deficiency, epimerase deficiency, and liver dysfunction also raise total galactose, so confirmatory enzyme and genetic testing are essential.
  \item (4) The Duarte variant produces reduced GALT activity (15--30\% of normal) but is benign and does not require dietary restriction.
\end{itemize}
\section*{Risks}
Minimal risk from heel stick or venipuncture: slight pain, bruising, or bleeding at the puncture site. Rarely: infection or, with improper heel-stick technique, local complications such as calcaneal injury.

\section*{What Happens After the Test}
Newborn screening results are typically available within a few days. A positive screen triggers immediate confirmatory quantitative GALT enzyme testing and, critically, institution of a galactose-free (soy) formula without waiting for confirmation. A confirmed diagnosis of classical galactosemia requires lifelong dietary galactose restriction, regular monitoring of red blood cell galactose-1-phosphate levels, and multidisciplinary follow-up (growth, development, speech, ovarian function in females, ophthalmology, and bone density), with genetic counseling for the family \cite{welling2017international}.

\section*{Understanding Results}

\subsection*{Normal Newborn Screen}
Normal GALT enzyme activity. Absence of elevated total galactose. Galactosemia is excluded. No further action.

\subsection*{Abnormal Newborn Screen (Positive Galactosemia Screen)}
Low or absent GALT enzyme activity in the dried blood spot. Elevated total galactose (galactose + galactose-1-phosphate). A positive screen requires immediate confirmatory testing (quantitative GALT enzyme activity in RBCs) and immediate clinical intervention. WHILE AWAITING CONFIRMATORY RESULTS, the infant must be placed immediately on a galactose-free formula (soy formula). Breast milk and cow milk formula MUST BE STOPPED immediately because continued galactose ingestion in a GALT-deficient infant causes hepatic failure, sepsis, and death within 1--2 weeks.

\subsection*{Classical Galactosemia (Confirmed Diagnosis)}
GALT enzyme activity in RBCs <1\% of normal (typically undetectable). Two pathogenic GALT mutations (e.g., Q188R/Q188R, Q188R/S135L). The infant must remain on a strict lifelong galactose-free diet: no breast milk, cow milk, cheese, yogurt, butter, organ meats, or other galactose-containing foods. Soy formula is the standard infant formula. Despite early treatment and excellent dietary compliance, many patients with classical galactosemia develop long-term complications: primary ovarian insufficiency (>80\% of females), developmental delay, speech dyspraxia (apraxia of speech), cognitive impairment, motor dysfunction (ataxia, tremor), and decreased bone mineral density. The pathogenesis of these complications is thought to involve endogenous galactose production (from UDP-glucose conversion in the body) that continues to produce low levels of galactose-1-phosphate accumulation despite dietary restriction \cite{berry2021classic}.

\subsection*{Duarte Variant (N314D/Duarte-2)}
Reduced GALT activity (approximately 15--30\% of normal). Elevated total galactose in the newborn period. The Duarte variant (compound heterozygosity of a classical mutation and the Duarte-2 allele) is a benign variant; patients are asymptomatic, do not develop clinical galactosemia, and do NOT require dietary restriction. Some states differentiate classical galactosemia from the Duarte variant on the newborn screen; others do not. Diagnostic testing by GALT genotyping is required.

\section*{Normal Results}
GALT enzyme activity in RBCs $\geq$10 U/g Hb (normal, $\geq$6--8 mcmol/min/g Hb). Galactose-1-phosphate in RBCs <1 mg/dL (87--130 mcg/g Hb). Total galactose in newborn blood spot <10 mg/dL. GALT genotyping: No pathogenic variants. Galactose tolerance: normal.

\section*{Follow-up Tests}
Quantitative GALT enzyme activity in RBCs (gold standard confirmatory test). GALT gene sequencing. Galactose-1-phosphate in RBCs for monitoring dietary compliance. Erythrocyte GALT isoform electrophoresis to distinguish variants. Urine reducing substances (not specific --- false positives and negatives; historical use only). Total galactose in plasma. Liver function tests (AST, ALT, PT, PTT), complete blood count (rule out sepsis), ophthalmologic examination (slit lamp exam for cataracts).

\section*{References}
\bibliographystyle{unsrtnat}
\bibliography{references}

\clearpage
\section*{Regulatory Status and Authorization Date}
This chapter describes a test category, not a specific commercial product. FDA, CDSCO, EU IVDR, and MHRA approval or registration dates therefore cannot be assigned without the assay manufacturer, product name, instrument, and intended use. For laboratory-developed tests, an individual agency approval date may not exist. See the Preface for the interpretation of regulatory status.

\clearpage
